% Partie : méthodes
% Chapitre : démonstrations
% Section : second exemple

\documentclass[../../../include/open-logic-section]{subfiles}

\begin{document}
\olfileid[fr]{mth}{prf}{ex2}

\olsection{Un autre exemple}

\begin{prop}
Si $A \subseteq C$, alors
$A \cup (C \setminus A) = C$.
\end{prop}

\begin{proof}
  Supposons $A \subseteq C$. Nous voulons
  montrer $A \cup (C \setminus A) = C$.
  \begin{quote}
    Observons d'abord que l'énoncé est une
    implication. Un quantificateur universel
    y est sous-entendu : la proposition vaut
    pour tous ensembles $A$ et $C$. Ces lettres
    désignent donc des ensembles arbitraires.
    Pour démontrer l'implication, nous supposons
    son antécédent et établissons son conséquent.

    Utilisons l'hypothèse $A \subseteq C$.
    Par définition de~$\subseteq$, cette inclusion signifie que tous
    les !!{element}s de~$A$ sont aussi !!{element}s de~$C$.
    Notons ce fait important, qui servira
    pendant toute la démonstration.
  \end{quote}
  Par définition de~$\subseteq$, puisque
  $A \subseteq C$, pour tout $z$, si
  $z \in A$, alors $z \in C$.
  \begin{quote}
    Nous avons déplié toutes les définitions
    présentes dans l'hypothèse. Passons à la
    conclusion $A \cup (C \setminus A) = C$.
    Comme dans l'exemple précédent, nous
    montrons que tout !!{element} de
    $A \cup (C \setminus A)$ est
    !!a{element} de~$C$ et,
    réciproquement, que tout !!{element}
    de~$C$ est !!a{element} de
    $A \cup (C \setminus A)$.
    Cela peut s'abréger en
    $A \cup (C \setminus A) \subseteq C$
    et $C \subseteq A \cup (C \setminus A)$.
    Nous faisons ici l'inverse du dépliage
    d'une définition pour rendre la
    démonstration plus facile à lire.
    Cette conjonction exige d'établir
    ses deux membres. Pour le premier,
    il faut montrer que tout !!{element}
    de $A \cup (C \setminus A)$ est
    !!a{element} de~$C$. Prenons donc
    un $z$ arbitraire appartenant à
    $A \cup (C \setminus A)$ et montrons
    qu'il appartient à~$C$. Par définition
    de $\cup$, nous pouvons conclure que
    $z \in A$ ou $z \in C \setminus A$.
    Vous commencez sans doute à voir
    comment procéder.
  \end{quote}
  $A \cup (C \setminus A) = C$ si et
  seulement si
  $A \cup (C \setminus A) \subseteq C$
  et $C \subseteq A \cup (C \setminus A)$.
  Montrons d'abord
  $A \cup (C \setminus A) \subseteq C$.
  Soit $z \in A \cup (C \setminus A)$.
  Alors $z \in A$ ou $z \in C \setminus A$.
  \begin{quote}
    Nous avons une disjonction et voulons
    en déduire $z \in C$. Distinguons
    donc deux cas.
  \end{quote}
  Cas 1 : $z \in A$. Comme nous savons
  que, pour tout $z$, $z \in A$ implique
  $z \in C$, nous obtenons $z \in C$.
  \begin{quote}
    Nous avons utilisé le fait tiré plus
    haut de l'hypothèse $A \subseteq C$.
    Le premier cas est achevé. Dans le
    second, $z \in C \setminus A$.
    Rappelons que $C \setminus A$ désigne
    la \emph{différence} des deux ensembles :
    l'ensemble des !!{element}s de~$C$
    qui ne sont pas !!{element}s de~$A$.
    Tout !!{element} de $C$ qui n'est
    pas dans~$A$ est, en particulier,
    !!a{element} de~$C$.
  \end{quote}
  Cas 2 : $z \in C \setminus A$. Cela
  signifie que $z \in C$ et
  $z \notin A$. En particulier,
  $z \in C$.
  \begin{quote}
    La première inclusion est démontrée.
    Passons à la seconde,
    $C \subseteq A \cup (C \setminus A)$.
    Supposons donc $z \in C$ et montrons
    que $z \in A \cup (C \setminus A)$.
  \end{quote}
  Soit maintenant $z \in C$. Nous
  devons montrer $z \in A$ ou
  $z \in C \setminus A$.
  \begin{quote}
    Puisque tous les !!{element}s de~$A$
    sont aussi !!{element}s de~$C$, et
    puisque $C \setminus A$ réunit les
    !!{element}s de~$C$ qui ne sont pas
    dans~$A$, il s'ensuit que $z$ est
    dans~$A$ ou dans $C \setminus A$.
    Si cette conclusion ne vous paraît
    pas déjà évidente, mieux vaut la
    démontrer étape par étape. Utilisons
    un fait simple que nous admettons :
    $z \in A$ ou $z \notin A$.
    C'est le «~principe du tiers exclu~» :
    pour tout énoncé~$p$, soit $p$ est
    vrai, soit sa négation est vraie.
    Ici, $p$ est l'énoncé $z \in A$.
    Cette disjonction nous permet de
    distinguer à nouveau deux cas.
  \end{quote}
  Soit $z \in A$, soit $z \notin A$.
  Dans le premier cas,
  $z \in A \cup (C \setminus A)$.
  Dans le second, $z \in C$ et
  $z \notin A$, donc
  $z \in C \setminus A$ et, par suite,
  $z \in A \cup (C \setminus A)$.
  \begin{quote}
    La démonstration est terminée :
    nous avons établi
    $A \cup (C \setminus A) = C$.
  \end{quote}
\end{proof}

\end{document}
